%! Author = WelJo
%! Date = 8/18/2026

\section{Fundamentals and Related Work}
\label{sec:related_works}

Existing methods for assessing CI resilience consider various behavioral aspects during disruptions~\cite{berdica2002,calvert2017,kaub2024,lickert2024}.
Resilience is not uniformly defined: Calvert and Snelder see it as a system's ability to cope with disruptions and restore functionality~\cite{calvert2017}, while Kaub et al.\ focus on road \textit{serviceability}~\cite{kaub2024}.
Bruneau et al.\ highlights its time-dependent nature, where performance decreases during disruptions and gradually recovers~\cite{bruneau2003}.
We follow this temporal perspective, evaluating functional losses and subsequent recovery.

\subsection{Interdependencies and Cascading Failures}
\label{subsec:interdependencies-and-cascading-failures}
Critical infrastructures are complex, interconnected networks~\cite{chiaradonna2007}.
Interdependencies are reciprocal relationships where one infrastructure's state influences another~\cite{laprie2007}.
These can be physical, cyber-related, or geographic~\cite{chiaradonna2007}.
Failures causing components of another infrastructure to fail are referred to as cascading failures~\cite{laprie2007}, which Buldyrev et al.\ show can propagate recursively~\cite{buldyrev2010}.
Lickert et al.\ investigate cascades triggered by power outages, noting that infrastructures may maintain functionality via emergency supply (e.g., hospital generators) for a limited duration, termed \textit{lifetimes}~\cite{lickert2024}.
Following these lifetimes and \textit{repair times}, systems transition between \textit{operative} and \textit{dysfunctional} states.

\subsection{Road Network Resilience and Accessibility}
\label{subsec:road-network-resilience-and-accessibility}
Emergency services depend on road network \textit{serviceability}, the availability of connections determining travel distances and costs~\cite{kaub2024}.
Disruptions reduce this serviceability, increasing travel times~\cite{berdica2002}.
While the \textit{Link Performance Index for Resilience} (LPIR) evaluates segments, it requires detailed data~\cite{calvert2017}.
Similarly, the \textit{Loss of Serviceability} (LoS) index evaluates resilience by comparing the shortest paths to analyze detours~\cite{kaub2024}.
However, LoS fails when nodes become isolated, as the distance becomes infinite and the index loses its explanatory value~\cite{kaub2024}.

Kaub et al.\ introduced the \textit{Robustness of Accessibility} (RoA) to evaluate entire urban networks~\cite{kaub2024}.
RoA compares the shortest path lengths in an intact network ($c(s_j,d_j)$) to a disrupted version ($c'(s'_j,d'_j)$). For each node pair, the ratio of these path lengths is weighted by $a_{\mathcal{C},j}$ and combined:

\begin{equation}
    \text{RoA}(\mathcal{C}) = \sum_{j=1}^{m} \left( a_{\mathcal{C},j} \cdot \frac{c(s_j,d_j)}{c'(s'_j, d'_j)} \right) \label{eq:equation}
\end{equation}

\noindent
RoA values range from 0 (complete failure) to 1 (full resilience).
If a destination becomes unreachable, the ratio is set to 0.
While Kaub et al.\ focus on road accessibility, they omit functional facility states.
Conversely, Lickert et al.\ model functional cascades without considering road accessibility~\cite{lickert2024}.
This paper bridges this gap by extending the RoA index with a time-dependent functional dimension, integrating CI cascading failures with the resulting impacts on urban accessibility.